\section{Data and Methods}
\label{sec:data-method}

\subsection{The r/AmItheAsshole and r/AITAH dataset}
\label{subsec:dataset}

We study two English-language moral-judgment communities on Reddit,
\emph{r/AmItheAsshole} (AITA) and \emph{r/AITAH} (AITAH). In both, an original
poster (OP) narrates an interpersonal conflict and asks whether they were in the
wrong; other users reply with a verdict drawn from the communities' shared tag
convention:
\emph{NTA} (Not the Asshole, exonerating the OP),
\emph{YTA} (You're the Asshole, blaming the OP),
\emph{NAH} (No Assholes Here),
\emph{ESH} (Everyone Sucks Here), and
\emph{INFO} (more information needed).

\paragraph{Verdict groups.}
Following community usage, we collapse the four scored tags into two moral
stances toward the OP: an \textbf{exculpating} group $\mathrm{NA}=\{\text{NTA},
\text{NAH}\}$ and an \textbf{inculpating} group $\mathrm{YA}=\{\text{YTA},
\text{ESH}\}$. The \emph{INFO} tag carries no judgment and is discarded
throughout. We restrict attention to \emph{top-level verdict comments}: direct
replies to the submission that carry one of the four scored tags.

\paragraph{Crowd verdict.}
For each post $p$ we summarize the community judgment by its \emph{YA share}
$\rho(p)\in[0,1]$, the fraction of its verdicts that blame the OP (i.e.\ that
fall in YA), together with the total number of verdicts the post received. To
ensure a stable crowd judgment, all analyses restrict to posts with at least
$200$ verdicts.

\paragraph{Scale.}
?? dataset size  -- number of posts, number of comments, number of pairs (top level comments and replies), number of users etc.

\paragraph{Certainty-annotated subset.}
For the language analyses we use, per community, an annotated subset of roughly
one thousand high-volume threads in which each top-level verdict comment carries
a human-annotated \emph{certainty} score on a four-point ordinal scale
($4$ = blunt and assertive, e.g.\ ``NTA. That's on her.''; $1$ = hedged and
conditional, e.g.\ ``YTA if you had the money, otherwise not.''). Because these
threads are large and mostly contested, we assign their consensus level
(\S\ref{subsec:method}) from the full crowd verdict rather than from the
annotated comments alone.

\paragraph{Cross-community structure.}
A minority of authors participate in both communities ($39{,}010$ users,
$4.34\%$ of the combined author set); among their posts, $8{,}025$ ($6.89\%$)
have a near-duplicate text counterpart in the other community. These
duplicate pairs---the same story judged by two different crowds---support the
community-level natural experiment analysed in \S[X].





\subsection{Method: four analytical dimensions}
\label{subsec:method}


\begin{table*}[t]
  \centering
  \begin{tabular}{@{}lll@{}}
    \toprule
    Dimension & Levels & Defined by \\
    \midrule
    (1) Alignment       & majority / minority        & comment verdict vs.\ dominant verdict \\
    (2) Consensus       & consensual / contested     & YA share ($\le.2$/$\ge.8$ vs.\ $.4$--$.6$) \\
    (3) Dissent direction & accusing / defending     & minority verdict $\times$ dominant verdict \\
    (4) Community       & AITA / AITAH               & source subreddit \\
    \bottomrule
  \end{tabular}
  \caption{The four analytical dimensions.}
  \label{tab:dimensions}
\end{table*}


Our analyses cross-cut the data along four binary dimensions. The first three
partition \emph{comments}, given the crowd verdict on their post; the fourth is
the \emph{community} itself, used as a replication axis.

\paragraph{Notation.}
For a post $p$ with YA share $\rho(p)$, define its \emph{dominant verdict}
\[
  D(p) =
  \begin{cases}
    \mathrm{YA} & \text{if } \rho(p) > 0.5,\\
    \mathrm{NA} & \text{if } \rho(p) < 0.5,
  \end{cases}
\]
dropping the measure-zero set of exactly tied posts ($\rho(p)=0.5$). For a
top-level verdict comment $c$ on post $p$, let $g(c)\in\{\mathrm{NA},\mathrm{YA}\}$
denote its verdict group.

\subsubsection*{(1) Majority vs.\ Minority --- alignment with the crowd}
A comment is \textbf{majority} if its verdict matches the crowd's dominant
verdict and \textbf{minority} (dissenting) otherwise:
\[
  \mathrm{side}(c) =
  \begin{cases}
    \text{majority} & \text{if } g(c) = D(p),\\
    \text{minority} & \text{if } g(c) \neq D(p).
  \end{cases}
\]
This is our central contrast: the minority is the voice that judges a given case
against the community's own verdict.

\subsubsection*{(2) Consensus: Consensual vs.\ Contested}
We bin posts by consensus strength, i.e.\ how lopsided the crowd verdict is,
using the YA share $\rho(p)$:
\begin{itemize}
  \item \textbf{Consensual} (strong consensus): $\rho(p)\le 0.20$
        (exculpating-dominant) or $\rho(p)\ge 0.80$ (inculpating-dominant).
  \item \textbf{Contested} (no consensus): $0.40\le\rho(p)\le 0.60$ (near tie).
\end{itemize}
The majority/minority split is only sharply meaningful under consensus; the
contested cohort serves as a control in which the two sides are near-symmetric.


\subsubsection*{(3) Accusing vs.\ Defending --- the direction of dissent}
Within the \textbf{minority}, dissent points in one of two opposite moral
directions relative to the crowd:
\begin{itemize}
  \item \textbf{Accusing}: $g(c)=\mathrm{YA}$ in an exculpating-dominant post
        --- the dissenter blames an OP the crowd exonerated.
  \item \textbf{Defending}: $g(c)=\mathrm{NA}$ in an inculpating-dominant post
        --- the dissenter exonerates an OP the crowd blamed.
\end{itemize}
The matching majority in each post type provides a same-context baseline, so any
accuse/defend difference is a property of the dissent itself, not of the post.

\subsubsection*{(4) AITA vs.\ AITAH --- community as a replication axis}
We run every analysis independently on the two communities and treat agreement
across them as evidence of a structural, rather than community-specific, effect.

\paragraph{The 18-hour rule.}
?? blahblah

\paragraph{Measures.}
For a top-level verdict comment we compute:
\begin{description}
  \item[Score] its net score, i.e.\ upvotes minus downvotes.
  \item[Direct replies] the number of comments that reply directly to it.
  \item[Reply depth] the number of nested reply levels beneath it
        ($0$ = no reply, $1$ = direct replies only, $2$ = a reply itself drew a
        reply, and so on), obtained by recursively traversing its reply subtree.
  \item[Thread score] the sum of net scores over the comment and its entire
        reply subtree---the total attention of the discussion it seeds.
  \item[Certainty] the annotated one-to-four assertiveness score (annotated
        subset only).
\end{description}
Because traversing the full reply forest is prohibitively expensive, the
depth- and thread-score measures are computed on a uniform random sample of
posts (the sample size is reported with each result); all other measures use the
full filtered data.
